\section{Ampelstatus}

\subsection{Ampel Report}
Ampelreport (prüfen!!):

\begin{table}[htbp]
    \centering
    \caption{Übersicht der der Ampelstatus} % Optional: Tabellenüberschrift
    \label{tab:ampelstatus}                  % Optional: Label für Verweise
    \begin{tabular}{lll}
        \toprule
        Status & relative Feuchte & Differenz absolute Feuchte \\
        \midrule
        Rot  & $\text{rF}_i \ge 65\,\% $       & $\Delta\text{aF} > 3 \frac{\text{g}}{\text{m}^3}$ \\ \addlinespace 
        Gelb & $65 > \text{rF}_i > 50$         & $\Delta\text{aF} \le 3\,\frac{\text{g}}{\text{m}^3}$ \\ \addlinespace
        Grün & $50 \ge \text{rF}_i$            & $\Delta\text{aF} \le 2\,\frac{\text{g}}{\text{m}^3}$ \\ \addlinespace
        nicht definiert & sonst                           & sonst \\  \bottomrule
    \end{tabular}
\end{table}


\subsection{Berechnung des Ampelstatus}
Bei der Berechnung des Ampelstatus wird der Absolutwert der Differenz der absoluten Feuchten in der Prozessluft $\text{aF}_p$ (process) und in der Raumluft $\text{aF}_a$ (ambient) gebildet.

\begin{align}
    \Delta\text{aF} 
    &=  \left| \text{aF}_p - \text{aF}_a \right| 
    \\[4ex]
    &=
    \frac{\text{rF}_p \cdot 611{,}2\, \frac{\text{J}}{\text{m}^3} \cdot e^{\dfrac{17{,}62 \cdot \vartheta_p}{\vartheta_p + 243{,}12}}}
    {461{,}52 \frac{\text{J}}{\text{kg} \cdot \text{K}} \cdot (\vartheta_p + 273{,}15)\, \text{K}}
    \cdot \frac{10}{1{,}204\, \frac{\text{kg}}{\text{m}^3}}
    - 
    \frac{\text{rF}_a \cdot 611{,}2\, \frac{\text{J}}{\text{m}^3} \cdot e^{\dfrac{17{,}62 \cdot \vartheta_a}{\vartheta_a + 243{,}12}}}
    {461{,}52 \frac{\text{J}}{\text{kg} \cdot \text{K}} \cdot (\vartheta_a + 273{,}15)\, \text{K}}
    \cdot \frac{10}{1{,}204\, \frac{\text{kg}}{\text{m}^3}}
    \\[4ex]
    &=
    \left[
    \text{rF}_p  \cdot
    \frac{ e^{\dfrac{17{,}62 \cdot \vartheta_p}{\vartheta_p + 243{,}12}}}
    { (\vartheta_p + 273{,}15)\, \text{K}}
    - 
    \text{rF}_a \cdot
    \frac{ e^{\dfrac{17{,}62 \cdot \vartheta_a}{\vartheta_a + 243{,}12}}}
    {(\vartheta_a + 273{,}15)\, \text{K}}
    \right]
    \cdot \frac{611{,}2\, \frac{\text{J}}{\text{m}^3}}{461{,}52 \frac{\text{J}}{\text{kg} \cdot \text{K}}}
    \cdot \frac{10}{1{,}204\, \frac{\text{kg}}{\text{m}^3}}
    \\[4ex]
     &=
    \left[
    \frac{\text{rF}_p}{\vartheta_a + 273{,}15} \cdot
    e^{f(\vartheta_p)}
    -
    \frac{\text{rF}_a}{\vartheta_p + 273{,}15} \cdot
    e^{f(\vartheta_a)}
    \right]
    \cdot \text{Konstante}   
\end{align}

\subsection{Graphische Darstellung}
\subsubsection{Ansatz 1}
2D Plot der Funktion im Intervall $[0 \dots 15]$:
\begin{equation}
    \Delta\text{aF} 
    =  \left| \text{aF}_p - \text{aF}_a \right|     
\end{equation}

